% Short running form; the full title is too wide for the head's right box.
\runninghead{The Church's Exodus toward the Judgment}
\properlane{last-exodus-of-the-church}{The Church's Exodus toward the Judgment: A People on the March Asks When}
\label{sec:last-exodus-of-the-church}

Heard from the Communion, the Mass asks a question of time. The singer has hoped
in God's salvation and in his word, and asks, \latin{quando facies de
persequentibus me iudicium?}, ``when wilt thou judge those who persecute me?''
The Epistle bids the Ephesians stand and resist \latin{in die malo}, girded about
the loins; the Gospel's king \latin{voluit rationem ponere}, willed to settle
accounts; the Alleluia sings Israel going out of Egypt; the Gradual confesses
the God who is \latin{a saeculo et usque in saeculum}; and the Postcommunion
names what has been eaten, \latin{immortalitatis alimonia}, the food of
immortality. St Augustine hears in the Communion's question the cry of the
martyrs in the Apocalypse, who ask how long until God judges, and are told to
wait until the number of their brethren is complete. St John Chrysostom reads
the girded loins of the Epistle as the girding with which Israel ate the
Passover on the night of the Exodus, and calls the present life ``nothing else
than a march and an encampment''. Read with them, the formulary is the song of a
people that has been brought out of Egypt and has not yet arrived. It asks when
the judgment will come, and it is answered that the judgment is the end toward
which it travels, and that until then it remains on earth, persecuted and helped,
fed on the way with the food of immortality.

\subsection{The question the Communion sings}

The Communion takes three half-verses from a strophe of Ps~118 whose speaker is
worn out with waiting. It leaves out his fainting, his failing eyes, his
shrivelled body and his near death, and it leaves out the first question of
verse 84, ``How many are the days of thy servant?'' What it sings is a hope, a
question and a cry: \latin{In salutari tuo anima mea, et in verbum tuum speravi:
quando facies de persequentibus me iudicium? iniqui persecuti sunt me, adiuva me,
Domine Deus meus}. The question is not how long the singer will live. It is when
God will judge. The difficulty is plain once the question is heard after the
Gospel. The parable has just told every hearer to forgive his brother from the
heart; the Communion asks for judgment on persecutors. Is the Church asking for
vengeance? And if the judgment is the last judgment, why ask when, when the Lord
told the disciples, ``It is not for you to know the times or moments'' (Acts 1:7)?

\subsection{The martyrs' plea and the long-suffering: St Augustine}

St Augustine reads the strophe as the voice of the whole people of God. The
singer who hopes in God's salvation is ``the chosen generation, the royal
priesthood, the holy nation, the peculiar people, longing for Christ from the
origin of the human race even unto the end of this world''; the first ages had
saints who desired his Incarnation, ``but these times, since He hath ascended
into heaven, have Saints who desire His manifestation to judge the quick and the
dead''. He says this of the soul's failing toward God's salvation, whose verb the
Communion omits; the salvation itself, \latin{in salutari tuo}, is sung. Of the
words the Communion does sing, \latin{in verbum tuum speravi}, he says: ``that
is, of Thy promise; a hope which causeth us to await with patience that which is
not seen by those who believe'' (\work{Enarrationes} 118, \latin{Caph}~81).

On verse 84 Augustine first names the voice: ``In the Apocalypse, these are the
words of the Martyrs, and long-suffering is enjoined them until the number of
their brethren be fulfilled.'' The martyrs under the altar cry, ``How long, O Lord
(Holy and True), dost thou not judge and revenge our blood on them that dwell on
the earth?'', and are told ``that they should rest for a little time till their
fellow servants and their brethren, who are to be slain even as they, should be
filled up'' (Apoc 6:10--11). The Communion's question is that cry, and its answer
is that answer: the judgment will come, and before it the Church waits. Augustine
then reads the two questions of the whole verse in order. ``The body of Christ
then is asking concerning its days, what they are to be in this world''; and
``when he had enquired concerning the days, he added also respecting the
judgment, showing indeed that the Church would exist on earth until the judgment,
when vengeance shall fall upon Her persecutors'' (118, \latin{Caph}~84). The first
of those two questions is the one the Communion leaves unsung, and Augustine's
``until the judgment'' comes from the order of the two. On the half that is sung,
his comment gives the martyrs' plea answered by long-suffering, and the vengeance
that falls at the judgment on the Church's persecutors.

On verse 86, whose second half the Communion sings, Augustine says that ``the
whole sense dependeth upon the foregoing'', and gives the cry for help its
object: ``help Thou me, that I may struggle for the truth even unto death''
(\latin{Caph}~86). The wicked persecuted ``nothing save the truth'', he says,
leaning on the first half of the verse, ``All thy statutes are truth'', which the
Communion leaves out.

\subsection{Not a curious question: Cassiodorus, St Robert Bellarmine and the Greek exposition}

Augustine is not alone in hearing the sung question as a desire for the end.
Cassiodorus says that the whole strophe is sung by \latin{peregrinus in hac terra
populus}, a people on pilgrimage in this land, confessing its desire for the
Lord's coming and the persecutions of the proud. On verse 84 the company of the
saints, knowing that the members of the Lord will be subject to the devil's
temptations until the end of the world, asks when this world will end, that the
vengeance prepared for the persecutors, \latin{id est diabolo cum ministris}, may
be completed; and he adds the sentence that settles the point, \latin{non quia
vitae suae tempus volebat agnoscere: sed celeritatem futuri iudicii postulabat}.
The singer was not asking to know the length of his own life; he was asking for
the swiftness of the judgment to come. He sets the question beside the Lord's
answer to the apostles that the times are not theirs to know (Acts 1:7), and
explains it by the blessed people's longing for the end of evils, which it could
see only after the end of this world (\work{Expositio Psalmorum} 118:84).

St Robert Bellarmine reads the sung half the same way: ``When will you judge and
condemn those that persecute and tempt me, that my persecution may cease on their
condemnation? He makes no curious inquiries as to the number of the days of his
life, or of those of the world, he merely expresses his desire of life
everlasting, when all the wicked \dots\ shall cease to harass the just''
(\work{Explanation of the Psalms}, Ps~118:84). On the question \latin{quando},
then, Augustine, Cassiodorus and Bellarmine agree: it asks for the judgment at the
end, and the hope that asks it waits with patience.

The exposition of the psalms printed under the name of St Athanasius does not
answer \latin{quando}, but it points the first and last sung verses toward the
same end. God's \latin{salutare}, it asks, what else can it be \latin{nisi
adventus Domini?}, and the word in which the singer hopes is the promise,
\latin{in promissionem} (PG~27, cols.~494--495). On the sung cry \latin{iniqui
persecuti sunt me} it says that the man persecuted unjustly, \latin{gnarus se in
divino tribunali mercedem accepturum}, knowing that he will receive his recompense
at the divine tribunal, \latin{gestit illud tempus videre}, longs to see that time
(col.~495). Its comment on verse 84 is said of the unsung question about the
servant's few days, and there it asks that Satan be crushed at once under his
feet, as the first reading notes beside Hilary. The exposition thus holds both:
the adversary crushed now, and the tribunal longed for at the end.

\subsection{Israel out of Egypt: the Alleluia and the Epistle's girded loins}

The Alleluia gives the waiting people a name and a road: \latin{In exitu Israel de
Aegypto}. Augustine teaches that ``no Christian consider himself alien to the
name of Israel''; that Egypt ``is often used for an emblem of this world; from
which we must spiritually withdraw''; and that ``as they did not depart thence,
until freed by Divine help; so no man is turned away in heart from this world,
unless aided by the gift of the Divine mercy''. That going out, he says, ``is
fulfilled in every faithful one in the daily travailings of the Church, in this
end of the world'' (\work{Enarrationes} 113.3--4). Cassiodorus takes the going out
as release from the bonds of sin and from the host of demons, and says that we
truly become Israel \latin{cum pompae mundi huius a nostra coeperint mente
discedere}, when the pomps of this world begin to leave our minds
(\work{Expositio Psalmorum} 113.1).

St John Chrysostom finds the same Exodus in the Epistle. On \latin{State ergo
succincti lumbos vestros in veritate} he asks why the Israelites ate the Passover
girded, with shoes on their feet and staves in their hands, and answers that the
guise was that of wayfarers: ``we were ready for our journey, we were just about
to go forth out of Egypt to the land of promise and we were ready for our
exodus''. ``This then is the historical type. But the reality is this; we too eat
a Passover, even Christ \dots\ We too ought to eat it, both sandalled and girded.
And why? That we too may be ready for our Exodus, for our departure hence.'' The
desert lies ahead, and the devil in it, ``to the end that those who have been
saved out of Egypt, those who have passed the Red Sea \dots\ may be taken and
destroyed by him''; but ``We have manna; yea rather not manna, but far more than
manna'', and ``the leader of this our Exodus is not Moses, but Jesus''
(\work{Homilies on Ephesians} 23). Chrysostom himself joins Paul's girded loins to
the Exodus and the Passover. And he draws the moral that governs this reading:
no one builds a splendid house in an encampment, and ``The present life is nothing
else than a march and an encampment'' (Hom.\ 23).

\subsection{The evil day, the eternal refuge and the king's reckoning}

The Epistle's \latin{dies malus} has two readings that this reading holds
together. Chrysostom says, ``By `evil day' he means the present life''
(Hom.\ 22). The continuation of \work{The Liturgical Year} grants that evil days
are many in every life and in the world's history, ``But, for every man, and for
the world at large, there is one evil day, evil beyond all the others: it is the
last day, the day of judgment'' (XI, p.~459). The two are one road: the whole
march is an evil day, and its end is the day that ends all others.

The Gradual names the God toward whom the march goes and in whom it shelters.
\latin{A saeculo et usque in saeculum tu es, Deus}: Augustine and Bellarmine both
mark the present tense, and Augustine calls that eternity the refuge to which we
fly from the mutability of time (\work{Enarrationes} 89.3; quoted under the
Gradual above); the exposition printed under St Athanasius's name finds in the
first verse the benefits God gave his people under Pharaoh, under Josue and in
Babylon, a refuge from generation to generation of Israel's own exodus and exile
(PG~27, col.~395). Fromage hears in the Gradual and the Alleluia together the
Church ``preparing the last exodus of the Church, which is leaving this world
turned infidel, and is hastening to the true land of promise'' (p.~463). The
psalm the Gradual begins goes on, in verses the Mass does not sing, to the
thousand years that are as yesterday and to the shortness of man's life; the
Gradual sings only the refuge and the eternity.

The Gospel's king gives the march its end. \latin{Qui voluit rationem ponere cum
servis suis}: St Thomas says that this reckoning begins when God first sends some
tribulation and ends at the day of judgment, \latin{Finis huius rationis erit in
die iudicii; principium quando primo inducit aliquam tribulationem} (\work{Super
Matthaeum} 18). Honorius Augustodunensis and William Durandus, expounding the
parable in their books, read the reckoning of the judgment (\work{Gemma animae}
IV.93; \work{Rationale} VI.139), and Fromage, at this Mass, sets beside the
Gospel the \latin{Dies irae}'s plea for forgiveness before the day of reckoning
(pp.~464--465), as the second reading shows. The reckoning the parable describes is the
judgment the Communion asks for, and the tribulations of the march are its
beginning.

\subsection{The rest of the Mass on the road}

The Introit is spoken by a people under sentence of death in a foreign empire,
far from the land, and in lines the Mass does not sing Mardochai reminds God of
``thy portion, which thou hast redeemed for thyself out of Egypt'' and asks ``that
we may live and praise thy name, O Lord'' (Esth 13:16--17). Fromage reads that
closing petition as the reason for the Collect's prayer for peace: the Church
asks to be free of adversity so that she may give God ``the united homage of
religion and good works'' (p.~456). The verse of the Introit blesses those who
walk, \latin{qui ambulant}, in the law of the Lord, and the Collect asks God to
keep his family on the way \latin{a cunctis adversitatibus}.

The Offertory's Job is the Church on the road. St Gregory reads Job's name,
``grieving'', as setting forth ``either our Mediator's Passion, or the travails of
Holy Church, which is harassed by the manifold toils of this present life''
(\work{Moralia}, Preface 16), and Schuster says that the history of Job ``is
repeated again and again throughout the history of the Church'' (p.~179). The
Secret offers the sacrifice by which God wills salvation to be restored,
\latin{salutem \dots\ restitui}: the salvation the Communion hopes in, \latin{in
salutari tuo}, is given on the way as well as at its end. And the Postcommunion
names the food: \latin{immortalitatis alimonia}. Chrysostom's ``we too eat a
Passover, even Christ'' and his ``far more than manna'' are said of the Epistle's
girded loins; the Postcommunion names the same food in its own words, the
nourishment of those who will not die, eaten by a people still on the march.
Bellarmine's gloss on the Communion's question, ``his desire of life
everlasting'', is the desire this food answers.

\subsection{The strongest difficulty}

Two objections meet this reading. The first is that it sets a cry for vengeance
after a Gospel of forgiveness. Augustine answers it from the martyrs themselves,
in the two passages the second reading quotes (\work{De sermone Domini in monte}
I.22.77; \work{Sermon} 58.8): the martyrs plead against the dominion of sin, not
against the men, and ``they might lawfully wish for the end of this world, in
which they had endured such martyrdoms''; and God's patience still waits for the
conversion of Christ's enemies. What this reading adds is the time. The
long-suffering enjoined on the martyrs is God's patience with their persecutors,
so the Church that asks when the judgment will come forgives while it waits,
because the delay is the time of the persecutors' conversion.

The second objection is that another Father reads the same verse otherwise, and
he does. St Hilary of Poitiers, whom the first reading follows, asks for a power
given in this life to crush the devil and his ministers, and says that this is not
the judgment of the resurrection, which he affirms: what the prophet now asks is
\latin{non illud utique quod in resurrectionis saeculo meminit futurum}
(\work{Tractatus in Psalmum} 118, \latin{Caph}~5, PL~9, cols.~573--574).
Augustine hears the martyrs' cry for that final judgment, ``when vengeance shall
fall upon Her persecutors'', the men who persecuted in the Church ``nothing save
the truth'' (\latin{Caph}~84, 86); Bellarmine, a plea that God ``judge and condemn
those that persecute and tempt me''; Cassiodorus joins Hilary's persecutors, the
devil with his ministers, to the judgment at the end of the world. The Fathers
differ in what is asked and in whom the plea names, not only in the time, which
each takes from the verse's unsung first question. Their answers are compatible
without being the same: Hilary affirms the judgment that Augustine and Bellarmine
await, and the men who persecute are, in the terms the first reading takes from
St Thomas, instruments moved by the devil. This reading follows the end-directed
answer, with Augustine's martyrs, Cassiodorus's swift judgment and the Athanasian
exposition's longing for the divine tribunal, without making it the only one.

\subsection{The four senses of this reading}

\begin{fourSenses}
\item[Literal.] Israel went out of Egypt, from a barbarous people. Mardochai prays
for a people under sentence of death. Paul tells the Ephesians to resist in the
evil day with their loins girt about with truth. The psalmist, hoping in God's
word, asks when God will judge his persecutors and begs his help. The king in the
parable wills to settle accounts with his servants.
\item[Allegorical.] The Exodus is the Church's going out of the world, led not by
Moses but by Jesus, with Christ for its Passover (Chrysostom); the Church whose
saints desire Christ's manifestation to judge asks when the judgment comes, as
the martyrs asked, and is told to wait (Augustine); the eternal God is its refuge
from the mutability of time (Augustine).
\item[Moral.] Live as on a march and in a camp, and build nothing in the desert
that you must leave (Chrysostom); withdraw in heart from Egypt, which no one does
without God's mercy (Augustine); forgive while God's patience waits for the
conversion of enemies (Augustine).
\item[Anagogical.] The judgment the Communion asks for, at which the king's
reckoning, begun in tribulation, is completed (Thomas); the land of promise toward
which the last exodus hastens (Fromage); the food of immortality fulfilled in the
life everlasting that Bellarmine hears the singer desire.
\end{fourSenses}
